%%%%%%%%%%%%%%%%%%%%%%%%%%%%%%%%%%%%%%%%%%%%%%%%%%%%%%%%%%%%%%%%%%%%%%%%%%%%%%%%
%%
\cleardoubleoddpage%  Make sure to start each chapter on a new odd page
\chapter{Methodology for Reproduction}
\label{sec:methodology}

\section{Selected Experiments from the Paper}

Both reproduced experiments are posed on the unit square $\Omega = (0,1)^2$
and share the same parametric structure: a coefficient or data field
depends on a parameter vector $x \in [-1,1]^d$, $d \in \{4,8\}$, through one
of two families used throughout the paper. The \emph{affine} family
expands the coefficient linearly in $x$, while the \emph{log-transformed}
family expands the logarithm of the coefficient, following the paper's
appendix~B formula
\[
    a(z,x) = \exp\Big(1 + x_1\Big(\tfrac{\sqrt{\pi\beta}}{2}\Big)^{1/2}
    + \sum_{j=2}^{d} \zeta_j\,\theta_j(z)\,x_j\Big),
\]
with $\zeta_j = (\sqrt{\pi\beta})^{1/2}\exp\!\big(-(\lfloor j/2\rfloor\pi\beta)^2/8\big)$,
$\beta_c = 1/8$, $\beta_p = \max(1, 2\beta_c)$, $\beta = \beta_c/\beta_p$, and
\[
    \theta_j(z) =
    \begin{cases}
        \sin\!\left(\left\lfloor \tfrac{j}{2} \right\rfloor \dfrac{\pi z_1}{\beta_p}\right), & j \text{ even}, \\[6pt]
        \cos\!\left(\left\lfloor \tfrac{j}{2} \right\rfloor \dfrac{\pi z_1}{\beta_p}\right), & j \text{ odd},
    \end{cases}
    \qquad j = 2, \dots, d.
\]
Each of the two PDEs is therefore instantiated in four cases (affine/log
$\times$ $d\in\{4,8\}$), giving eight cases in total across the two
experiments.

\paragraph{Parametric elliptic diffusion.} The diffusion problem is posed
in mixed form: find $(\sigma, u) \in H(\mathrm{div};\Omega) \times L^2(\Omega)$
such that
\[
    \int_\Omega \frac{\sigma \cdot \tau}{a(x)}\,\mathrm{d}\Omega
    + \int_\Omega u\,\mathrm{div}(\tau)\,\mathrm{d}\Omega
    = G(\tau)
    \quad \forall\, \tau \in H(\mathrm{div};\Omega),
\]
\[
    \int_\Omega v\,\mathrm{div}(\sigma)\,\mathrm{d}\Omega
    = -\int_\Omega f v\,\mathrm{d}\Omega
    \quad \forall\, v \in L^2(\Omega),
\]
with forcing $f = 10$ and Dirichlet data $u = 0.5$ on the bottom edge of
$\Omega$ and $u = 0$ on the remaining boundary. In this mixed formulation
the Dirichlet data enters \emph{naturally}, through the boundary functional
$G(\tau) = \int_{\partial\Omega} g\,(\tau\cdot n)\,\mathrm{d}s$, rather than
as an essential boundary condition on a primal $H^1(\Omega)$ space. This is
a structural difference from a standard finite-difference or primal-FEM
discretization, and it has to be respected for the reproduction to match
the paper's actual formulation.

\paragraph{Navier--Stokes--Brinkman.} The NSB system is posed in the
four-field mixed form of Gatica, N\'{u}\~{n}ez and Ruiz-Baier
\cite{gatica2022nsb}, the paper's own cited source for this formulation:
find $(u, t, \sigma, \gamma) \in \mathbf{L}^4(\Omega) \times L^2_{\mathrm{tr}}(\Omega)
\times H_0(\mathrm{div}_{4/3};\Omega) \times L^2_{\mathrm{skew}}(\Omega)$
satisfying the nonlinear residual system with convective term
$\int (u\otimes u):s$, viscosity coupling
$\lambda \int a(x)\,t : s$, and a Brinkman reaction term $\eta\, u$, where
$\lambda = 0.1$ and $\eta(x) = 10 + z_1^2 + z_2^2$. A Nitsche-type penalty
term $\kappa \langle (\sigma + u\otimes u)n, \tau n\rangle_{\partial\Omega_{\mathrm{out}}} = 0$
with $\kappa = 10^4$ enforces the outlet condition weakly, forcing is
$f = (0,-1)$, and an inlet velocity profile is prescribed on the top
boundary. Because the primary unknown $u$ is sought in
$\mathbf{L}^4(\Omega)$, this experiment is Banach-valued in the sense of
Section~\ref{sec:background}, unlike the Hilbert-valued diffusion problem.

\section{Data Generation / Dataset Construction}

\paragraph{Finite element discretization.} Both PDEs are discretized with
FEniCS~2019.1.0 (legacy \texttt{dolfin}), run inside a WSL2 Ubuntu
environment since the required legacy FEniCS build is not available on
native Windows. The paper's appendix gives the continuous weak-form
equations for both problems, but that alone does not fix a computational
realization: it does not say which mesh to use, what polynomial degree
each finite element field should have, or (for one ambiguous printed
formula, the NSB inlet velocity) what literal numeric value to use. For
this reason, this reproduction uses the mesh and finite element spaces
exactly as implemented in the paper authors' own released code, instead of
building an independent mesh and guessing at a calibrated element pair.
Both problems are solved on the authors' own mesh, which has 143 vertices
and 244 triangular cells and was obtained directly from the supervisor and
loaded without modification. The diffusion problem uses the same element
pair as their code, $BDM_2 \times DG_1$ \cite{brezzidouglasmarini1985}
(Brezzi--Douglas--Marini degree 2 for the pseudostress $\sigma$, degree-1
discontinuous Galerkin for $u$), which gives 2622 total degrees of freedom
(1890 for $\sigma$, 732 for $u$) and matches the paper's own reported
figure exactly. The NSB problem uses the Arnold--Falk--Winther-type
element family \cite{arnoldfalkwinther2007} from \cite{gatica2022nsb} at
the same polynomial degree as the authors' code: $BDM_2$ per pseudostress
row, $DG_1$-vector for velocity $u$, $DG_1$ for skew vorticity $\gamma$,
and $DG_2$-vector (3 components) for the trace-free strain $t$. This gives
1464 degrees of freedom for $u$ and 244 for the pressure $p$ (projected
into $DG_0$), again matching the paper's own reported figures exactly. The
inlet velocity is the authors' constant $(0.1, 0.0)$ profile, and the
pressure gauge is fixed exactly as in their code, a $DG_0$ projection with
no additional mean shift. The nonlinear system is solved with Newton's
method via \texttt{NonlinearVariationalSolver}.

\paragraph{Sparse-grid test protocol.} Rather than drawing an independent
random test sample, the test error is evaluated on a fixed, deterministic
Clenshaw--Curtis sparse grid \cite{clenshawcurtis1960} over $[-1,1]^d$,
built with the Tasmanian toolkit. Sparse grids construct a
quadrature/interpolation point set for a $d$-dimensional hypercube from
one-dimensional rules via Smolyak's combination technique
\cite{smolyak1963}, which grows far more slowly with $d$ than a full
tensor-product grid. That is the property that makes a deterministic,
exhaustive test discretization feasible at $d=8$ at all. The same
sparse-grid test set is shared across all training-set sizes $m$ and all
twelve random trials within a case, exactly matching the paper's protocol
of reusing one deterministic test discretization per case. The sparse-grid
level is not stated numerically in the paper's own text, but the authors'
own batch scripts (\texttt{Pbatch.sh}, \texttt{Nbatch\_u.sh}) fix it
explicitly, and this reproduction now matches their choice exactly: level
5 (1105 points) at $d=4$, level 4 (3937 points) at $d=8$, for both
experiments.

\paragraph{Bochner-norm test error.} The relative test error is measured in
the FEM-mass-matrix-weighted norm $\lVert \cdot \rVert_Y$ induced by the
function space's mass matrix $M$ ($\lVert v \rVert_Y^2 = v^\top M v$),
combined with the sparse-grid quadrature weights $w_i$, rather than a
naive Euclidean sum-of-squares over raw degree-of-freedom vectors averaged
uniformly over the test points. This matters in practice for two reasons.
Both meshes are non-uniform, so an unweighted sum-of-squares would
over-weight degrees of freedom on smaller elements. And the Smolyak
combination technique used to build the sparse grid produces genuinely
negative quadrature weights at a real fraction of the test points (720 out
of 3937 at $d=8$, 336 out of 1105 at $d=4$, in this reproduction's own
grids), so a poorly-fit model can push the accumulated weighted sum
negative. The paper authors' own code guards against this by taking the
absolute value of the accumulated sum on each side of the ratio before
the square root, and this reproduction does the same; without that guard,
the square root of a negative number silently produces NaN instead of a
large but finite error. The mass matrix is assembled once per case during
data generation and stored as a sidecar file so that the Windows-side
training code never needs a FEniCS installation.

\paragraph{Dataset scale.} Each of the eight cases (2 PDEs $\times$ 2
coefficient families $\times$ 2 dimensions) provides training data for
fourteen training-set sizes $m$ and twelve independent random trials per
size, plus one shared deterministic sparse-grid test set, matching the
paper's full experimental matrix.

\section{Model Architecture}

Following the paper's encoder--network--decoder structure
(Section~\ref{sec:background}), each surrogate is a fully connected
network mapping the parameter vector $x \in [-1,1]^d$ directly to the
finite element degrees of freedom of the discretized solution. Six
architecture/activation combinations are used, matching the paper: two
depth/width configurations (4 hidden layers of width 40, and 10 hidden
layers of width 100) combined with three activation functions (ReLU, ELU,
and tanh). Weights are initialized with He/Kaiming-uniform initialization
\cite{he2015delving} using a ReLU-derived bound for all three activations,
matching Keras' \texttt{HeUniform} initializer, which is activation-agnostic.
This required an explicit fix in both the PyTorch and JAX implementations,
whose default Kaiming/Xavier conventions are activation-aware and would
otherwise silently diverge from the paper's stated initialization
(Sections~\ref{sec:pytorchreimplementation} and~\ref{sec:jaximplementation}).

\section{Training Procedure}

Both frameworks train with the Adam optimizer \cite{kingma2015adam} and an
exponentially decaying learning rate schedule, minimizing the mean squared
error between predicted and reference finite element degrees of freedom.
For the diffusion experiment, training runs for up to 60{,}000 epochs, with
an early stop once the training loss falls below a tolerance of
$5 \times 10^{-7}$, exactly as stated in the paper's appendix. For the NSB
experiment, this budget is reduced to 20{,}000 epochs and a
$5 \times 10^{-6}$ tolerance, a documented deviation
(Appendix~\ref{app:an-appendix}) made after direct measurement showed NSB
runs routinely needing tens of thousands of epochs under the same decaying
schedule to approach the paper's tolerance, which made the full paper-scale
matrix intractable on the hardware available for this practical work.

The paper's Appendix~A.2(iv) describes a checkpoint/restore rule in prose
that reads, on a first pass, like it has two triggers: save whenever
either the current loss is the best seen so far, or the ratio between the
current loss and the loss at the last checkpoint drops below $1/8$.
Cross-checking the authors' own training callback shows this is not
actually the case. The ratio only controls how often an expensive
test-error metric gets logged during training, and has no effect on which
weights are saved. The rule that actually governs checkpointing and
restoration is a single trigger: save whenever the current loss is the
best seen so far, and after training terminates, restore the checkpointed
weights if the final-epoch loss is worse than the checkpoint's loss. This
reproduction implements exactly that single-trigger rule, identically for
PyTorch and JAX, so that any residual difference between the two
frameworks' results reflects their numerical/optimizer behavior rather
than a difference in stopping policy.

\section{Reproducibility Controls}

Every one of the twelve trials per case uses a distinct, explicitly
recorded random seed that controls both the training-set subsample drawn
for that trial and the corresponding network initialization; no trial
reuses another trial's seed. The sparse-grid test set, by contrast, is
fully deterministic and identical across all trials and training-set
sizes within a case, so that observed variation in test error reflects
only the training-set sample and initialization, not test-set sampling
noise. Aggregation across trials follows the paper's convention of
geometric mean and geometric standard deviation
($\exp(\mathrm{mean}(\log x))$, $\exp(\mathrm{std}(\log x))$) rather than
arithmetic statistics, since relative errors are naturally log-scaled
quantities and are plotted on log-log axes.
%%
%%%%%%%%%%%%%%%%%%%%%%%%%%%%%%%%%%%%%%%%%%%%%%%%%%%%%%%%%%%%%%%%%%%%%%%%%%%%%%%%
